    \begin{center}
        {\large \textbf{Вариант №22 \\
        }}
    \end{center}

    \textbf{\underline{Цель работы:} изучение способов задания инфологической модели данных и создания структуры базы данных в заданной предметной области.
    } 

    \begin{enumerate}
        \item \hyperref[sec:1]{Выполнить анализ предметной области, выделить основные сущности, атрибуты и связи.}
        \item \hyperref[sec:2]{Создать диаграмму «сущность — связь» в нотации Чена.}
        \item \hyperref[sec:3]{Самостоятельно изучить нотацию IDEF1X для представления диаграммы «сущность-связь». Создать схему базы данных в нотации IDEF1X.}
        \item \hyperref[sec:4]{Разработать структуру базы данных и составить описание столбцов таблиц базы данных, включающее: имя столбца, назначение (какие данные хранятся), тип данных, допускает ли столбец пустые значения.}
    \end{enumerate}